%
%============= CHAPTER BREAK =================%
%
\chapter{Aelius Aristides (ca 147--181 CE)}%
\vspace*{3mm} %
\justifying %

\noindent{}Born in rural Bithynia in 117 CE, and granted citizenship as a child by Emperor Hadrian, %
a grant made also to his father, %
Aelius Aristides eventually became one of the most famous orators in the empire. %
As a young man he lived and learned in Smyrna. %
After travels and further education in Alexandria, Egypt, %
he planned to travel west, %
but ill health prevented. %
Multiple years of seeking healing and dwelling in an Asklepion led to a modicum of recovery %
and a variety of religious experiences. %
Around 154 CE, he found health sufficient to travel, %
stayed in Athens and then in Rome, %
where he also took on students. \\[1.5mm] %

\noindent{}\textbf{περιαστράπτω} / \textit{peristraptō} / “to flash around” in the TLG %
has its results conflated with another lemma, περιστρέφω / \textit{peristrephō} / “to crown around”. %
Once we filter out hits of the latter, %
the first non-Christian use of the former is found in the enormously famous %
\textit{Panathenaic Oration} of Aelius Aristides. %
Scholars are somewhat divided about when the oration was composed: %
Trapp maintained 155 CE, but Oliver placed it in 167 CE, %
and Follet in 168 CE.\endnote{Oudot 2008:32n7, with further literature.} %
This “encomium to the city of Athens, phrased as if delivered as part of the celebration %
of the four-yearly Great Panathenaea”,\endnote{Trapp 2017:3.} %
was not likely performed at the official ceremonies, %
not least because of its overwhelming length.\endnote{\textit{Ibid}., 4.} %
Athenian cultural superiority echoes across the entire speech, %
but Athensʼ military achievements take up its bulk (§§ 25--330).\endnote{\textit{Ibid}., 5.} %

In section 165 of the oration, %
Aristides commemorates the victory of the Athenian Greeks at the Battle of Salamis, %
noting just before (§160) that the Spartans and other Greeks were not present at the battle. %
It is in this episode, a dramatic, decisive, yet still early turning point in the Greek wars with Persia, %
where the new lemma occurs. %

\vspace*{2.5mm} 
 \begin{quote}κατὰ οὖν τοὺς ἐκπίπτοντας ἐκ τῆς ναυμαχίας ἐτάχθη τὰ κύρια τῆς στρατιᾶς, %
 ἄνδρες Περσῶν οἱ πρῶτοι, %
 καὶ κατέσχον τὴν νῆσον ἐφεστηκότες ὡς ἐφ’ ἑτοίμοις. %     
 ὡς δὲ τῶν πρώτων ἀπέτυχον
 καὶ μετὰ τῶν Ἑλλήνων ἦν ἡ τύχη καὶ \textbf{περιέστραπτο} ἡ πεῖρα, %
 ἀνὴρ Ἀθηναίων εἷς ἐθελοντὴς τὸν κίνδυνον ὑφίστατο %
 καὶ λαβὼν τοὺς ὄντας αὐτῶν ἐν Σαλαμῖνι τοὺς πρεσβύτας %
 ἀποβιβάσας εἰς τὴν νῆσον κτείνει πᾶν τοῦτο τὸ Περσικόν. (\textit{Panathen}. 165)\endnote{%
 Behr :158.} \\[1mm] %
 Therefore, against those cast ashore from the naval battle, %
 the lords of the army were stationed ---  %
 the foremost men of the Persians ---, %
 and they were holding down the island, 
 having taken a stand as over prepared matters. %
 Now as those of the front lines failed, %
 and Fortune was with the Greeks, %
 and the travail \textbf{was flashing around} them, %
 one man of the Athenians, a volunteer, %
 undertook the danger, %
 and taking those of them in Salamis who were the elders, %
 and disembarking them on the island, he slew this entire Persicon. %
 \end{quote} % 
 \vspace*{2.5mm}

\noindent{}In Aristidesʼ telling, this turn of the battle tide %
led Xerxes to behold Athens with amazement, then make a hasty retreat (§166). %

The same new lemma is used in the famous, first telling of Saulʼs encounter with the risen Jesus %
on the road to Damascus (Acts 9:3--4):\endnote{%
The lemma is repeated in the second narration of Saulʼs Damascus Road encounter at Acts 22:6: %
περιαστράψαι / \textit{periastrapsai}.} 

\vspace*{2.5mm} 
 \begin{quote}ἐν δὲ τῷ πορεύεσθαι ἐγένετο αὐτὸν ἐγγίζειν τῇ Δαμασκῷ, %
 ἐξαίφνης τε αὐτὸν \textbf{περιήστραψεν} φῶς ἐκ τοῦ οὐρανοῦ, %
 ⁴ καὶ πεσὼν ἐπὶ τὴν γῆν ἤκουσεν φωνὴν λέγουσαν αὐτῷ, %
 Σαοὺλ Σαούλ, τί με διώκεις; \clearpage %
 Now as he was going, it happened as he drew near to Damascus, %
 suddenly light from the heaven \textbf{flashed around} him, %
 ⁴ and falling upon the ground, he heard a voice saying to him, %
 Saoul, Saoul, why are you pursuing me? %
 \end{quote} % 
 \vspace*{2.5mm}

\noindent{}In previous scholarship, van der Horst wrote an entire monograph comparing the themes %
and vocabulary of Aelius Aristides to the New Testament %
and in regard to Acts 9:3 noted a parallel of light as a divine manifestation, %
but failed to notice this unique lemma connection.\endnote{%
van der Horst 1980:34.} %
The imitation, if strategic and deliberate, %
suggests that the Christ who conquered Saul on the Road to Damascus %
is not merely divine, but is also a personification of the power of Greek culture. %
Christ overcame Saul, flashing all around him, %
the same way that the Athenian Greeks overcame Xerxes. %
Saul subsequently could only marvel at Christ near Damascus, %
just as Xerxes could only stand in astonishment at the miraculous power of Athens Herself. %
It could also suggest that Damascus is implicitly depicted in Acts %
as a center of Hellenistic culture and political power. %

The editor(s) of 4 Maccabees --- which may well be the same person(s) as the editor(s) of Acts ---
also likely made use of Aristidesʼ imagery and vocabulary. %
As I described in \textit{Apostle Paul Atreides} (chp 8): %

\vspace*{2mm} 
 \begin{quote}Apollonius of Tarsus, introduced as governor of Coelesyria and Phoenicia (2Mac 3:5), or of Syria, Phoenicia, and Cilicia (4Mac 4:2), is worth considering as dual theophoric and mimetic antecedent to Paul. Both share Tarsus as home city. Apollonius is implicated in an attempt by Simon, a displaced Judean high-priest, to summon King Seleucis to plunder and thus profane the Jerusalem temple, similar to the temple-focused accusations Paul faces in Acts both in Jerusalem and Ephesus. In Maccabean lore this Apollonius also fell to the ground due to a divine epiphany involving horses! %
 In 2 Maccabees (3:24--30), Heliodorus (not to be confused with the author of the Greek novel several centuries later) represented Seleucis and encountered an epiphanic divine army counter-assault on horseback. %
 But in 4 Maccabees, Apollonius of Tarsus now takes the role as emissary, %
 potential plunderer of Jerusalemʼs temple, and victim of divine, epiphanic downfall.
\end{quote} % 
 \vspace*{2mm}

\noindent{}The relevant quotation follows: %

\vspace*{3mm}
\begin{quote} %
\textls[-12]{ἀνιόντος τε μετὰ καθωπλισμένης τῆς στρατιᾶς τοῦ Ἀπολλωνίου πρὸς τὴν τῶν χρημάτων %
ἁρπαγὴν οὐρανόθεν ἔφιπποι προυφάνησαν ἄγγελοι \textbf{περιαστράπτοντες} τοῖς ὅπλοις καὶ πολὺν αὐτοῖς %
φόβον τε καὶ τρόμον ἐνιέντες. ¹¹ καταπεσών γέ τοι ἡμιθανὴς ὁ Ἀπολλώνιος ἐπὶ τὸν πάμφυλον %
τοῦ ἱεροῦ περίβολον τὰς χεῖρας ἐξέτεινεν εἰς τὸν οὐρανὸν καὶ μετὰ δακρύων τοὺς Εβραίους %
παρεκάλει ὅπως περὶ αὐτοῦ προσευξάμενοι τὸν οὐράνιον ἐξευμενίσωνται στρατόν.} (4Mac 4:10--11) \\[1mm] %
\textls[-10]{And when Apollonius with his well-armed battalion was going up for the seizure %
of the financial property, from heaven horse-riding angels appeared ahead, %
\textbf{flashing around} with their weapons and infusing into them much fear and trembling. %
¹¹ Then Apollonius, falling down half-dead upon the all-tribes courtyard of the temple, %
stretched out hands to the heaven, and with tears was begging the Hebrews that they, praying for him, %
would propitiate the heavenly army.} %
\end{quote} % 
\vspace*{3mm}

Here Apollonius (and the Seleucid Greek ruler Antiochus IV Epiphanes, whom Apollonius represents) %
become a second Xerxes, and Jerusalem becomes a second Athens. %
The audacious endeavor to pillage the Jerusalem temple, %
a wealthy religious center akin to Athens, %
is met, stopped, and undone with divine flashes of lightning. %